\chapter{Security Analysis}

\section{Security of AES-CTR-HMAC with RISC Zero Proofs}

\subsection{Scope}

This chapter gives a game-based proof sketch, not a full formal reduction. The
goal is to explain why the public transcript produced by the implemented
pipeline does not reveal which of two equal-length plaintexts was encrypted in a
single challenge encryption, assuming the usual cryptographic and zkVM
assumptions hold. In standard terminology, the game below is an
IND-1-CPA-style game: it captures one challenge encryption and does not give the
adversary a general encryption oracle.

The public transcript is written as
\((\mathrm{aad}, n, c, t, y, \pi)\), where \(\mathrm{aad}\) is associated data or
context, \(n\) is the AES-CTR IV/nonce, \(c\) is ciphertext, \(t\) is the
truncated HMAC tag, \(y\) is public output metadata, and \(\pi\) is the RISC Zero
receipt. This matches the implementation framing
\texttt{aad || iv || ciphertext}; the notation \(n\) is used for the IV/nonce in
the proof sketch. For authentication, the AAD is length-delimited before it is
concatenated with the nonce and ciphertext, so the MAC input is parsed
unambiguously.

The confidentiality claim requires \(y\) to be independent of the challenge bit,
except for public metadata that is already allowed to leak, such as message
length. If a future guest journals plaintext-dependent application output, that
output must be modeled separately because it can reveal information regardless of
the encryption layer.

This chapter also states the integrity and execution-binding assumptions used by
the report, but it does not prove full chosen-ciphertext security, sender
identity, replay protection, or side-channel resistance.

\subsection{Assumptions}

The proof sketch relies on the following assumptions:

\begin{itemize}
\item \textbf{Nonce discipline:} AES-CTR IV/nonce values are unique under each
encryption key. Equivalently, the same AES counter block is never encrypted twice
under the same key, including across multi-block messages and counter increments.
If two messages reuse the same counter stream, XORing their ciphertexts cancels
the keystream and leaks the XOR of the plaintexts.
\item \textbf{AES-192-CTR confidentiality:} AES-192 is treated as a secure block
cipher/PRF, so CTR mode is IND-CPA secure when nonces are never reused
\cite{nistfips197,nistsp80038a,katz2020introduction}.
\item \textbf{HMAC security:} HMAC-SHA256, truncated to 192 bits in this
implementation, is treated as a secure PRF for tag indistinguishability and as a
secure MAC for tamper-detection reasoning \cite{rfc2104,nistfips1981}.
\item \textbf{Key separation:} encryption and MAC keys are distinct.
\item \textbf{zkVM proof properties:} receipt verification is sound for the
expected guest image ID, and the proof does not leak non-journaled witness data
beyond the public journal and verification result.
\end{itemize}

The 192-bit tag truncation gives a generic blind-forgery term of approximately
\(q_v / 2^{192}\) for \(q_v\) verification attempts, before accounting for any
advantage against HMAC itself. This is negligible for the scale considered in
this dissertation, but it is still a stated assumption rather than an empirical
benchmark result.

\subsection{Security Game}

The one-challenge confidentiality game is:

\begin{enumerate}
\item The adversary \(\mathcal{A}\) outputs two equal-length messages
\(m_0, m_1\) and public associated data \(\mathrm{aad}\).
\item The challenger samples \(b \leftarrow \{0,1\}\), a fresh nonce \(n\), and
secret keys \(k_E, k_M\).
\item The guest relation computes:
\begin{align*}
c &= \mathrm{AES\text{-}CTR}_{k_E}(m_b; n) \\
\ell_{\mathrm{aad}} &= \mathrm{EncodeLen}(|\mathrm{aad}|) \\
t &= \mathrm{Trunc}_{192}(\mathrm{HMAC}_{k_M}(\ell_{\mathrm{aad}} \parallel \mathrm{aad} \parallel n \parallel c)).
\end{align*}
\item The guest commits public outputs \((\mathrm{aad}, n, c, t, y)\), where
\(y\) is the journaled public output of the computation, such as public metadata
or an application result explicitly intended to be public. For this
confidentiality claim, \(y\) must be independent of \(b\) except through allowed
leakage such as message length. The prover emits a receipt \(\pi\) for the
expected guest image.
\item \(\mathcal{A}\) receives \((\mathrm{aad}, n, c, t, y, \pi)\) and outputs a
guess \(b'\).
\end{enumerate}

This is not the full multi-query IND-CPA game. A full IND-CPA definition would
also give \(\mathcal{A}\) access to an encryption oracle before and after the
challenge, with nonce uniqueness enforced across all oracle and challenge
encryptions. The proof sketch below should therefore be read as covering the
one-encryption case used by the transcript-generation argument.

The length encoding \(\ell_{\mathrm{aad}}\) is not needed for the one-challenge
confidentiality argument, but it is needed for later ciphertext-integrity
reasoning. Without an unambiguous boundary between \(\mathrm{aad}\), \(n\), and
\(c\), an adversary could potentially reinterpret a valid tuple by moving the
nonce and first ciphertext block across field boundaries, for example using
\(\mathrm{aad}' = \mathrm{aad} \parallel n\), \(n' = c_0\), and
\(c' = c_1 \parallel \cdots\). Length-delimited AAD prevents this splice from
preserving the authenticated byte string.

The advantage is:

\[
\mathrm{Adv}_{\mathcal{A}} :=
\left| \Pr[b' = b] - \tfrac{1}{2} \right|.
\]

\subsection{Game Sequence}

\paragraph{Game \(G_0\) (real execution).}
This is the implemented transcript generation described above.

\paragraph{Game \(G_1\) (replace AES-CTR keystream).}
Let \(\mathrm{Ctr}(n,i)\) denote the \(i\)-th AES input block derived from nonce
\(n\) and counter value \(i\). In the real game, AES-CTR encrypts the sequence of
counter blocks to produce the keystream

\[
S = \mathrm{AES}_{k_E}(\mathrm{Ctr}(n,0)) \parallel
\mathrm{AES}_{k_E}(\mathrm{Ctr}(n,1)) \parallel \cdots,
\]

truncated to \(|m_b|\) bits. Game \(G_1\) replaces this AES-generated keystream
with a uniformly random string \(r\) of the same length:

\[
c = m_b \oplus r, \quad r \leftarrow \{0,1\}^{|m_b|}.
\]

Under the AES-CTR PRF assumption and nonce uniqueness:

\[
|\Pr[G_0] - \Pr[G_1]| \leq
\mathrm{Adv}_{\mathrm{PRF}}^{\mathrm{AES}}(\mathcal{A}).
\]

Because \(m_0\) and \(m_1\) have equal length, \(c\) is now distributed
independently of \(b\).

\paragraph{Game \(G_2\) (replace HMAC tag).}
Replace the HMAC tag on
\(\ell_{\mathrm{aad}} \parallel \mathrm{aad} \parallel n \parallel c\) with a
random 192-bit value. The tag is already a function of public values after
\(G_1\), so this hop mainly records the assumption that the truncated HMAC output
does not introduce a distinguishable signal about \(b\):

\[
|\Pr[G_1] - \Pr[G_2]| \leq
\mathrm{Adv}_{\mathrm{PRF}}^{\mathrm{HMAC}}(\mathcal{A}).
\]

For integrity rather than confidentiality, the corresponding forgery term is
bounded informally by the HMAC unforgeability advantage plus the truncation term
\(q_v / 2^{192}\).

\paragraph{Game \(G_3\) (simulate receipt).}
Let \(\mathrm{Sim}\) be the zero-knowledge simulator for the proof system. It
takes only the public journal values and produces a simulated receipt
indistinguishable from a real receipt to any efficient verifier. Game \(G_3\)
replaces the real receipt \(\pi\) with a simulated receipt \(\tilde{\pi}\) over
the same public journal values:

\[
\tilde{\pi} \leftarrow \mathrm{Sim}(\mathrm{aad}, n, c, t, y).
\]

By the zero-knowledge property for non-journaled witness data:

\[
|\Pr[G_2] - \Pr[G_3]| \leq
\mathrm{Adv}_{\mathrm{ZK}}(\mathcal{A}).
\]

This step does not hide the journal; it only says the receipt should not reveal
extra information about secret inputs beyond the values intentionally committed
to the journal.

\paragraph{Game \(G_4\) (ideal).}
In the final game, \((c, t, y, \pi)\) is independent of \(b\) under the stated
condition that \(y\) does not encode plaintext-dependent information. Therefore:

\[
\Pr[G_4 : b' = b] = \tfrac{1}{2}.
\]

\subsection{Final Bound}

By the triangle inequality, the one-challenge confidentiality advantage is
bounded by:

\begin{align*}
\mathrm{Adv}_{\mathcal{A}}
&\leq \mathrm{Adv}_{\mathrm{PRF}}^{\mathrm{AES}}
+ \mathrm{Adv}_{\mathrm{PRF}}^{\mathrm{HMAC}}
+ \mathrm{Adv}_{\mathrm{ZK}}.
\end{align*}

For transcript integrity and execution binding, the report relies on separate
terms: HMAC unforgeability plus \(q_v / 2^{192}\) for truncated-tag forgery, and
RISC Zero soundness plus image-ID verification for the claim that the expected
guest image generated or checked the journaled tuple.

\subsection{Conclusion}

Under nonce-respecting AES-CTR, separated AES and HMAC keys, PRF/MAC security of
HMAC-SHA256-192, and RISC Zero proof soundness and non-journal leakage
assumptions, the public transcript does not reveal which equal-length challenge
plaintext was encrypted in the IND-1-CPA game. The result is deliberately
scoped: it supports the report's confidentiality and transcript-binding claims
for the implemented benchmark pipeline, but it is not a full multi-query
IND-CPA proof, a full AEAD proof, or a deployment-level network security proof.

\hfill \(\rule{0.8ex}{0.8ex}\)
